% -*- compile-command: "rubber -d graphs.tex" -*-
\documentclass{tufte-handout}

\usepackage{algo-activity}

\title{\thecourse: Introduction to Graphs}
\date{}

\begin{document}

\maketitle

\begin{model}{Graphs and graph terms}{graphs}
  \begin{center}
  \begin{minipage}{0.45\textwidth}
    \begin{pgfpicture}
\pgfpathrectangle{\pgfpointorigin}{\pgfqpoint{150.0000bp}{150.0000bp}}
\pgfusepath{use as bounding box}
\begin{pgfscope}
\pgfsetlinewidth{0.5655bp}
\definecolor{sc}{rgb}{0.0000,0.0000,0.0000}
\pgfsetstrokecolor{sc}
\pgfsetmiterjoin
\pgfsetbuttcap
\pgfpathqmoveto{110.8756bp}{99.7024bp}
\pgfpathqlineto{104.4652bp}{71.6757bp}
\pgfusepathqstroke
\end{pgfscope}
\begin{pgfscope}
\pgfsetlinewidth{0.5655bp}
\definecolor{sc}{rgb}{0.0000,0.0000,0.0000}
\pgfsetstrokecolor{sc}
\pgfsetmiterjoin
\pgfsetbuttcap
\pgfpathqmoveto{110.8756bp}{99.7024bp}
\pgfpathqlineto{82.2135bp}{95.6883bp}
\pgfusepathqstroke
\end{pgfscope}
\begin{pgfscope}
\pgfsetlinewidth{0.5655bp}
\definecolor{sc}{rgb}{0.0000,0.0000,0.0000}
\pgfsetstrokecolor{sc}
\pgfsetmiterjoin
\pgfsetbuttcap
\pgfpathqmoveto{104.4652bp}{71.6757bp}
\pgfpathqlineto{82.2135bp}{95.6883bp}
\pgfusepathqstroke
\end{pgfscope}
\begin{pgfscope}
\pgfsetlinewidth{0.5655bp}
\definecolor{sc}{rgb}{0.0000,0.0000,0.0000}
\pgfsetstrokecolor{sc}
\pgfsetmiterjoin
\pgfsetbuttcap
\pgfpathqmoveto{104.4652bp}{71.6757bp}
\pgfpathqlineto{117.3643bp}{40.4385bp}
\pgfusepathqstroke
\end{pgfscope}
\begin{pgfscope}
\pgfsetlinewidth{0.5655bp}
\definecolor{sc}{rgb}{0.0000,0.0000,0.0000}
\pgfsetstrokecolor{sc}
\pgfsetmiterjoin
\pgfsetbuttcap
\pgfpathqmoveto{117.3643bp}{40.4385bp}
\pgfpathqlineto{103.4723bp}{13.4302bp}
\pgfusepathqstroke
\end{pgfscope}
\begin{pgfscope}
\pgfsetlinewidth{0.5655bp}
\definecolor{sc}{rgb}{0.0000,0.0000,0.0000}
\pgfsetstrokecolor{sc}
\pgfsetmiterjoin
\pgfsetbuttcap
\pgfpathqmoveto{82.2135bp}{95.6883bp}
\pgfpathqlineto{51.3110bp}{110.7828bp}
\pgfusepathqstroke
\end{pgfscope}
\begin{pgfscope}
\pgfsetlinewidth{0.5655bp}
\definecolor{sc}{rgb}{0.0000,0.0000,0.0000}
\pgfsetstrokecolor{sc}
\pgfsetmiterjoin
\pgfsetbuttcap
\pgfpathqmoveto{51.3110bp}{110.7828bp}
\pgfpathqlineto{33.3304bp}{136.2671bp}
\pgfusepathqstroke
\end{pgfscope}
\begin{pgfscope}
\pgfsetlinewidth{0.5655bp}
\definecolor{sc}{rgb}{0.0000,0.0000,0.0000}
\pgfsetstrokecolor{sc}
\pgfsetmiterjoin
\pgfsetbuttcap
\pgfpathqmoveto{33.3304bp}{136.2671bp}
\pgfpathqlineto{5.3272bp}{128.3841bp}
\pgfusepathqstroke
\end{pgfscope}
\begin{pgfscope}
\pgfsetlinewidth{0.5655bp}
\definecolor{sc}{rgb}{0.0000,0.0000,0.0000}
\pgfsetstrokecolor{sc}
\pgfsetmiterjoin
\pgfsetbuttcap
\pgfpathqmoveto{5.3272bp}{128.3841bp}
\pgfpathqlineto{21.0578bp}{104.1213bp}
\pgfusepathqstroke
\end{pgfscope}
\begin{pgfscope}
\pgfsetlinewidth{0.5655bp}
\definecolor{sc}{rgb}{0.0000,0.0000,0.0000}
\pgfsetstrokecolor{sc}
\pgfsetmiterjoin
\pgfsetbuttcap
\pgfpathqmoveto{21.0578bp}{104.1213bp}
\pgfpathqlineto{51.3110bp}{110.7828bp}
\pgfusepathqstroke
\end{pgfscope}
\begin{pgfscope}
\pgfsetlinewidth{0.5655bp}
\definecolor{sc}{rgb}{0.0000,0.0000,0.0000}
\pgfsetstrokecolor{sc}
\pgfsetmiterjoin
\pgfsetbuttcap
\pgfpathqmoveto{117.3643bp}{40.4385bp}
\pgfpathqlineto{144.8617bp}{27.3170bp}
\pgfusepathqstroke
\end{pgfscope}
\begin{pgfscope}
\definecolor{fc}{rgb}{0.0000,0.0000,0.0000}
\pgfsetfillcolor{fc}
\pgftransformshift{\pgfqpoint{150.0000bp}{24.8650bp}}
\pgftransformscale{1.0000}
\pgftext[]{$9$}
\end{pgfscope}
\begin{pgfscope}
\definecolor{fc}{rgb}{0.0000,0.0000,0.0000}
\pgfsetfillcolor{fc}
\pgfsetlinewidth{0.5655bp}
\definecolor{sc}{rgb}{0.0000,0.0000,0.0000}
\pgfsetstrokecolor{sc}
\pgfsetmiterjoin
\pgfsetbuttcap
\pgfpathqmoveto{146.0004bp}{27.3170bp}
\pgfpathqcurveto{146.0004bp}{27.9458bp}{145.4906bp}{28.4556bp}{144.8617bp}{28.4556bp}
\pgfpathqcurveto{144.2328bp}{28.4556bp}{143.7230bp}{27.9458bp}{143.7230bp}{27.3170bp}
\pgfpathqcurveto{143.7230bp}{26.6881bp}{144.2328bp}{26.1783bp}{144.8617bp}{26.1783bp}
\pgfpathqcurveto{145.4906bp}{26.1783bp}{146.0004bp}{26.6881bp}{146.0004bp}{27.3170bp}
\pgfpathclose
\pgfusepathqfillstroke
\end{pgfscope}
\begin{pgfscope}
\definecolor{fc}{rgb}{0.0000,0.0000,0.0000}
\pgfsetfillcolor{fc}
\pgftransformshift{\pgfqpoint{18.8963bp}{98.8542bp}}
\pgftransformscale{1.0000}
\pgftext[]{$8$}
\end{pgfscope}
\begin{pgfscope}
\definecolor{fc}{rgb}{0.0000,0.0000,0.0000}
\pgfsetfillcolor{fc}
\pgfsetlinewidth{0.5655bp}
\definecolor{sc}{rgb}{0.0000,0.0000,0.0000}
\pgfsetstrokecolor{sc}
\pgfsetmiterjoin
\pgfsetbuttcap
\pgfpathqmoveto{22.1965bp}{104.1213bp}
\pgfpathqcurveto{22.1965bp}{104.7501bp}{21.6867bp}{105.2599bp}{21.0578bp}{105.2599bp}
\pgfpathqcurveto{20.4289bp}{105.2599bp}{19.9191bp}{104.7501bp}{19.9191bp}{104.1213bp}
\pgfpathqcurveto{19.9191bp}{103.4924bp}{20.4289bp}{102.9826bp}{21.0578bp}{102.9826bp}
\pgfpathqcurveto{21.6867bp}{102.9826bp}{22.1965bp}{103.4924bp}{22.1965bp}{104.1213bp}
\pgfpathclose
\pgfusepathqfillstroke
\end{pgfscope}
\begin{pgfscope}
\definecolor{fc}{rgb}{0.0000,0.0000,0.0000}
\pgfsetfillcolor{fc}
\pgftransformshift{\pgfqpoint{0.0000bp}{130.3929bp}}
\pgftransformscale{1.0000}
\pgftext[]{$7$}
\end{pgfscope}
\begin{pgfscope}
\definecolor{fc}{rgb}{0.0000,0.0000,0.0000}
\pgfsetfillcolor{fc}
\pgfsetlinewidth{0.5655bp}
\definecolor{sc}{rgb}{0.0000,0.0000,0.0000}
\pgfsetstrokecolor{sc}
\pgfsetmiterjoin
\pgfsetbuttcap
\pgfpathqmoveto{6.4659bp}{128.3841bp}
\pgfpathqcurveto{6.4659bp}{129.0130bp}{5.9561bp}{129.5228bp}{5.3272bp}{129.5228bp}
\pgfpathqcurveto{4.6983bp}{129.5228bp}{4.1885bp}{129.0130bp}{4.1885bp}{128.3841bp}
\pgfpathqcurveto{4.1885bp}{127.7552bp}{4.6983bp}{127.2454bp}{5.3272bp}{127.2454bp}
\pgfpathqcurveto{5.9561bp}{127.2454bp}{6.4659bp}{127.7552bp}{6.4659bp}{128.3841bp}
\pgfpathclose
\pgfusepathqfillstroke
\end{pgfscope}
\begin{pgfscope}
\definecolor{fc}{rgb}{0.0000,0.0000,0.0000}
\pgfsetfillcolor{fc}
\pgftransformshift{\pgfqpoint{35.2343bp}{141.6327bp}}
\pgftransformscale{1.0000}
\pgftext[]{$6$}
\end{pgfscope}
\begin{pgfscope}
\definecolor{fc}{rgb}{0.0000,0.0000,0.0000}
\pgfsetfillcolor{fc}
\pgfsetlinewidth{0.5655bp}
\definecolor{sc}{rgb}{0.0000,0.0000,0.0000}
\pgfsetstrokecolor{sc}
\pgfsetmiterjoin
\pgfsetbuttcap
\pgfpathqmoveto{34.4690bp}{136.2671bp}
\pgfpathqcurveto{34.4690bp}{136.8960bp}{33.9592bp}{137.4058bp}{33.3304bp}{137.4058bp}
\pgfpathqcurveto{32.7015bp}{137.4058bp}{32.1917bp}{136.8960bp}{32.1917bp}{136.2671bp}
\pgfpathqcurveto{32.1917bp}{135.6382bp}{32.7015bp}{135.1284bp}{33.3304bp}{135.1284bp}
\pgfpathqcurveto{33.9592bp}{135.1284bp}{34.4690bp}{135.6382bp}{34.4690bp}{136.2671bp}
\pgfpathclose
\pgfusepathqfillstroke
\end{pgfscope}
\begin{pgfscope}
\definecolor{fc}{rgb}{0.0000,0.0000,0.0000}
\pgfsetfillcolor{fc}
\pgftransformshift{\pgfqpoint{55.0020bp}{115.1176bp}}
\pgftransformscale{1.0000}
\pgftext[]{$5$}
\end{pgfscope}
\begin{pgfscope}
\definecolor{fc}{rgb}{0.0000,0.0000,0.0000}
\pgfsetfillcolor{fc}
\pgfsetlinewidth{0.5655bp}
\definecolor{sc}{rgb}{0.0000,0.0000,0.0000}
\pgfsetstrokecolor{sc}
\pgfsetmiterjoin
\pgfsetbuttcap
\pgfpathqmoveto{52.4496bp}{110.7828bp}
\pgfpathqcurveto{52.4496bp}{111.4117bp}{51.9398bp}{111.9215bp}{51.3110bp}{111.9215bp}
\pgfpathqcurveto{50.6821bp}{111.9215bp}{50.1723bp}{111.4117bp}{50.1723bp}{110.7828bp}
\pgfpathqcurveto{50.1723bp}{110.1539bp}{50.6821bp}{109.6441bp}{51.3110bp}{109.6441bp}
\pgfpathqcurveto{51.9398bp}{109.6441bp}{52.4496bp}{110.1539bp}{52.4496bp}{110.7828bp}
\pgfpathclose
\pgfusepathqfillstroke
\end{pgfscope}
\begin{pgfscope}
\definecolor{fc}{rgb}{0.0000,0.0000,0.0000}
\pgfsetfillcolor{fc}
\pgftransformshift{\pgfqpoint{100.8682bp}{8.3673bp}}
\pgftransformscale{1.0000}
\pgftext[]{$4$}
\end{pgfscope}
\begin{pgfscope}
\definecolor{fc}{rgb}{0.0000,0.0000,0.0000}
\pgfsetfillcolor{fc}
\pgfsetlinewidth{0.5655bp}
\definecolor{sc}{rgb}{0.0000,0.0000,0.0000}
\pgfsetstrokecolor{sc}
\pgfsetmiterjoin
\pgfsetbuttcap
\pgfpathqmoveto{104.6110bp}{13.4302bp}
\pgfpathqcurveto{104.6110bp}{14.0591bp}{104.1012bp}{14.5689bp}{103.4723bp}{14.5689bp}
\pgfpathqcurveto{102.8434bp}{14.5689bp}{102.3336bp}{14.0591bp}{102.3336bp}{13.4302bp}
\pgfpathqcurveto{102.3336bp}{12.8013bp}{102.8434bp}{12.2915bp}{103.4723bp}{12.2915bp}
\pgfpathqcurveto{104.1012bp}{12.2915bp}{104.6110bp}{12.8013bp}{104.6110bp}{13.4302bp}
\pgfpathclose
\pgfusepathqfillstroke
\end{pgfscope}
\begin{pgfscope}
\definecolor{fc}{rgb}{0.0000,0.0000,0.0000}
\pgfsetfillcolor{fc}
\pgftransformshift{\pgfqpoint{121.4966bp}{44.3550bp}}
\pgftransformscale{1.0000}
\pgftext[]{$3$}
\end{pgfscope}
\begin{pgfscope}
\definecolor{fc}{rgb}{0.0000,0.0000,0.0000}
\pgfsetfillcolor{fc}
\pgfsetlinewidth{0.5655bp}
\definecolor{sc}{rgb}{0.0000,0.0000,0.0000}
\pgfsetstrokecolor{sc}
\pgfsetmiterjoin
\pgfsetbuttcap
\pgfpathqmoveto{118.5030bp}{40.4385bp}
\pgfpathqcurveto{118.5030bp}{41.0674bp}{117.9932bp}{41.5772bp}{117.3643bp}{41.5772bp}
\pgfpathqcurveto{116.7354bp}{41.5772bp}{116.2256bp}{41.0674bp}{116.2256bp}{40.4385bp}
\pgfpathqcurveto{116.2256bp}{39.8096bp}{116.7354bp}{39.2998bp}{117.3643bp}{39.2998bp}
\pgfpathqcurveto{117.9932bp}{39.2998bp}{118.5030bp}{39.8096bp}{118.5030bp}{40.4385bp}
\pgfpathclose
\pgfusepathqfillstroke
\end{pgfscope}
\begin{pgfscope}
\definecolor{fc}{rgb}{0.0000,0.0000,0.0000}
\pgfsetfillcolor{fc}
\pgftransformshift{\pgfqpoint{78.8185bp}{91.1179bp}}
\pgftransformscale{1.0000}
\pgftext[]{$2$}
\end{pgfscope}
\begin{pgfscope}
\definecolor{fc}{rgb}{0.0000,0.0000,0.0000}
\pgfsetfillcolor{fc}
\pgfsetlinewidth{0.5655bp}
\definecolor{sc}{rgb}{0.0000,0.0000,0.0000}
\pgfsetstrokecolor{sc}
\pgfsetmiterjoin
\pgfsetbuttcap
\pgfpathqmoveto{83.3522bp}{95.6883bp}
\pgfpathqcurveto{83.3522bp}{96.3171bp}{82.8424bp}{96.8269bp}{82.2135bp}{96.8269bp}
\pgfpathqcurveto{81.5846bp}{96.8269bp}{81.0748bp}{96.3171bp}{81.0748bp}{95.6883bp}
\pgfpathqcurveto{81.0748bp}{95.0594bp}{81.5846bp}{94.5496bp}{82.2135bp}{94.5496bp}
\pgfpathqcurveto{82.8424bp}{94.5496bp}{83.3522bp}{95.0594bp}{83.3522bp}{95.6883bp}
\pgfpathclose
\pgfusepathqfillstroke
\end{pgfscope}
\begin{pgfscope}
\definecolor{fc}{rgb}{0.0000,0.0000,0.0000}
\pgfsetfillcolor{fc}
\pgftransformshift{\pgfqpoint{99.6704bp}{68.6059bp}}
\pgftransformscale{1.0000}
\pgftext[]{$1$}
\end{pgfscope}
\begin{pgfscope}
\definecolor{fc}{rgb}{0.0000,0.0000,0.0000}
\pgfsetfillcolor{fc}
\pgfsetlinewidth{0.5655bp}
\definecolor{sc}{rgb}{0.0000,0.0000,0.0000}
\pgfsetstrokecolor{sc}
\pgfsetmiterjoin
\pgfsetbuttcap
\pgfpathqmoveto{105.6039bp}{71.6757bp}
\pgfpathqcurveto{105.6039bp}{72.3046bp}{105.0941bp}{72.8144bp}{104.4652bp}{72.8144bp}
\pgfpathqcurveto{103.8364bp}{72.8144bp}{103.3266bp}{72.3046bp}{103.3266bp}{71.6757bp}
\pgfpathqcurveto{103.3266bp}{71.0469bp}{103.8364bp}{70.5371bp}{104.4652bp}{70.5371bp}
\pgfpathqcurveto{105.0941bp}{70.5371bp}{105.6039bp}{71.0469bp}{105.6039bp}{71.6757bp}
\pgfpathclose
\pgfusepathqfillstroke
\end{pgfscope}
\begin{pgfscope}
\definecolor{fc}{rgb}{0.0000,0.0000,0.0000}
\pgfsetfillcolor{fc}
\pgftransformshift{\pgfqpoint{115.0702bp}{103.5521bp}}
\pgftransformscale{1.0000}
\pgftext[]{$0$}
\end{pgfscope}
\begin{pgfscope}
\definecolor{fc}{rgb}{0.0000,0.0000,0.0000}
\pgfsetfillcolor{fc}
\pgfsetlinewidth{0.5655bp}
\definecolor{sc}{rgb}{0.0000,0.0000,0.0000}
\pgfsetstrokecolor{sc}
\pgfsetmiterjoin
\pgfsetbuttcap
\pgfpathqmoveto{112.0143bp}{99.7024bp}
\pgfpathqcurveto{112.0143bp}{100.3313bp}{111.5045bp}{100.8411bp}{110.8756bp}{100.8411bp}
\pgfpathqcurveto{110.2467bp}{100.8411bp}{109.7369bp}{100.3313bp}{109.7369bp}{99.7024bp}
\pgfpathqcurveto{109.7369bp}{99.0736bp}{110.2467bp}{98.5638bp}{110.8756bp}{98.5638bp}
\pgfpathqcurveto{111.5045bp}{98.5638bp}{112.0143bp}{99.0736bp}{112.0143bp}{99.7024bp}
\pgfpathclose
\pgfusepathqfillstroke
\end{pgfscope}
\end{pgfpicture}

  \end{minipage}
  \begin{minipage}{0.45\textwidth}
    \begin{pgfpicture}
\pgfpathrectangle{\pgfpointorigin}{\pgfqpoint{150.0000bp}{150.0000bp}}
\pgfusepath{use as bounding box}
\begin{pgfscope}
\pgfsetlinewidth{0.5987bp}
\definecolor{sc}{rgb}{0.0000,0.0000,0.0000}
\pgfsetstrokecolor{sc}
\pgfsetmiterjoin
\pgfsetbuttcap
\pgfpathqmoveto{144.0766bp}{67.2085bp}
\pgfpathqlineto{112.0079bp}{88.2694bp}
\pgfusepathqstroke
\end{pgfscope}
\begin{pgfscope}
\pgfsetlinewidth{0.5987bp}
\definecolor{sc}{rgb}{0.0000,0.0000,0.0000}
\pgfsetstrokecolor{sc}
\pgfsetmiterjoin
\pgfsetbuttcap
\pgfpathqmoveto{137.4218bp}{117.1461bp}
\pgfpathqlineto{112.0079bp}{88.2694bp}
\pgfusepathqstroke
\end{pgfscope}
\begin{pgfscope}
\pgfsetlinewidth{0.5987bp}
\definecolor{sc}{rgb}{0.0000,0.0000,0.0000}
\pgfsetstrokecolor{sc}
\pgfsetmiterjoin
\pgfsetbuttcap
\pgfpathqmoveto{112.0079bp}{88.2694bp}
\pgfpathqlineto{68.6020bp}{80.1602bp}
\pgfusepathqstroke
\end{pgfscope}
\begin{pgfscope}
\pgfsetlinewidth{0.5987bp}
\definecolor{sc}{rgb}{0.0000,0.0000,0.0000}
\pgfsetstrokecolor{sc}
\pgfsetmiterjoin
\pgfsetbuttcap
\pgfpathqmoveto{68.6020bp}{80.1602bp}
\pgfpathqlineto{45.7604bp}{115.0036bp}
\pgfusepathqstroke
\end{pgfscope}
\begin{pgfscope}
\pgfsetlinewidth{0.5987bp}
\definecolor{sc}{rgb}{0.0000,0.0000,0.0000}
\pgfsetstrokecolor{sc}
\pgfsetmiterjoin
\pgfsetbuttcap
\pgfpathqmoveto{45.7604bp}{115.0036bp}
\pgfpathqlineto{21.8889bp}{144.1840bp}
\pgfusepathqstroke
\end{pgfscope}
\begin{pgfscope}
\pgfsetlinewidth{0.5987bp}
\definecolor{sc}{rgb}{0.0000,0.0000,0.0000}
\pgfsetstrokecolor{sc}
\pgfsetmiterjoin
\pgfsetbuttcap
\pgfpathqmoveto{68.6020bp}{80.1602bp}
\pgfpathqlineto{45.1884bp}{40.9871bp}
\pgfusepathqstroke
\end{pgfscope}
\begin{pgfscope}
\pgfsetlinewidth{0.5987bp}
\definecolor{sc}{rgb}{0.0000,0.0000,0.0000}
\pgfsetstrokecolor{sc}
\pgfsetmiterjoin
\pgfsetbuttcap
\pgfpathqmoveto{45.1884bp}{40.9871bp}
\pgfpathqlineto{6.9699bp}{48.0099bp}
\pgfusepathqstroke
\end{pgfscope}
\begin{pgfscope}
\pgfsetlinewidth{0.5987bp}
\definecolor{sc}{rgb}{0.0000,0.0000,0.0000}
\pgfsetstrokecolor{sc}
\pgfsetmiterjoin
\pgfsetbuttcap
\pgfpathqmoveto{45.1884bp}{40.9871bp}
\pgfpathqlineto{25.5620bp}{6.4905bp}
\pgfusepathqstroke
\end{pgfscope}
\begin{pgfscope}
\pgfsetlinewidth{0.5987bp}
\definecolor{sc}{rgb}{0.0000,0.0000,0.0000}
\pgfsetstrokecolor{sc}
\pgfsetmiterjoin
\pgfsetbuttcap
\pgfpathqmoveto{45.1884bp}{40.9871bp}
\pgfpathqlineto{70.6877bp}{11.9140bp}
\pgfusepathqstroke
\end{pgfscope}
\begin{pgfscope}
\definecolor{fc}{rgb}{0.0000,0.0000,0.0000}
\pgfsetfillcolor{fc}
\pgftransformcm{1.0000}{0.0000}{0.0000}{1.0000}{\pgfqpoint{75.3605bp}{6.5863bp}}
\pgftransformscale{1.0000}
\pgftext[]{j}
\end{pgfscope}
\begin{pgfscope}
\definecolor{fc}{rgb}{0.0000,0.0000,0.0000}
\pgfsetfillcolor{fc}
\pgfsetlinewidth{0.5987bp}
\definecolor{sc}{rgb}{0.0000,0.0000,0.0000}
\pgfsetstrokecolor{sc}
\pgfsetmiterjoin
\pgfsetbuttcap
\pgfpathqmoveto{72.1050bp}{11.9140bp}
\pgfpathqcurveto{72.1050bp}{12.6968bp}{71.4705bp}{13.3313bp}{70.6877bp}{13.3313bp}
\pgfpathqcurveto{69.9050bp}{13.3313bp}{69.2704bp}{12.6968bp}{69.2704bp}{11.9140bp}
\pgfpathqcurveto{69.2704bp}{11.1312bp}{69.9050bp}{10.4967bp}{70.6877bp}{10.4967bp}
\pgfpathqcurveto{71.4705bp}{10.4967bp}{72.1050bp}{11.1312bp}{72.1050bp}{11.9140bp}
\pgfpathclose
\pgfusepathqfillstroke
\end{pgfscope}
\begin{pgfscope}
\definecolor{fc}{rgb}{0.0000,0.0000,0.0000}
\pgfsetfillcolor{fc}
\pgftransformcm{1.0000}{0.0000}{0.0000}{1.0000}{\pgfqpoint{22.0576bp}{0.3310bp}}
\pgftransformscale{1.0000}
\pgftext[]{i}
\end{pgfscope}
\begin{pgfscope}
\definecolor{fc}{rgb}{0.0000,0.0000,0.0000}
\pgfsetfillcolor{fc}
\pgfsetlinewidth{0.5987bp}
\definecolor{sc}{rgb}{0.0000,0.0000,0.0000}
\pgfsetstrokecolor{sc}
\pgfsetmiterjoin
\pgfsetbuttcap
\pgfpathqmoveto{26.9793bp}{6.4905bp}
\pgfpathqcurveto{26.9793bp}{7.2733bp}{26.3447bp}{7.9078bp}{25.5620bp}{7.9078bp}
\pgfpathqcurveto{24.7792bp}{7.9078bp}{24.1446bp}{7.2733bp}{24.1446bp}{6.4905bp}
\pgfpathqcurveto{24.1446bp}{5.7077bp}{24.7792bp}{5.0732bp}{25.5620bp}{5.0732bp}
\pgfpathqcurveto{26.3447bp}{5.0732bp}{26.9793bp}{5.7077bp}{26.9793bp}{6.4905bp}
\pgfpathclose
\pgfusepathqfillstroke
\end{pgfscope}
\begin{pgfscope}
\definecolor{fc}{rgb}{0.0000,0.0000,0.0000}
\pgfsetfillcolor{fc}
\pgftransformcm{1.0000}{0.0000}{0.0000}{1.0000}{\pgfqpoint{0.0000bp}{49.2906bp}}
\pgftransformscale{1.0000}
\pgftext[]{h}
\end{pgfscope}
\begin{pgfscope}
\definecolor{fc}{rgb}{0.0000,0.0000,0.0000}
\pgfsetfillcolor{fc}
\pgfsetlinewidth{0.5987bp}
\definecolor{sc}{rgb}{0.0000,0.0000,0.0000}
\pgfsetstrokecolor{sc}
\pgfsetmiterjoin
\pgfsetbuttcap
\pgfpathqmoveto{8.3872bp}{48.0099bp}
\pgfpathqcurveto{8.3872bp}{48.7926bp}{7.7527bp}{49.4272bp}{6.9699bp}{49.4272bp}
\pgfpathqcurveto{6.1871bp}{49.4272bp}{5.5526bp}{48.7926bp}{5.5526bp}{48.0099bp}
\pgfpathqcurveto{5.5526bp}{47.2271bp}{6.1871bp}{46.5925bp}{6.9699bp}{46.5925bp}
\pgfpathqcurveto{7.7527bp}{46.5925bp}{8.3872bp}{47.2271bp}{8.3872bp}{48.0099bp}
\pgfpathclose
\pgfusepathqfillstroke
\end{pgfscope}
\begin{pgfscope}
\definecolor{fc}{rgb}{0.0000,0.0000,0.0000}
\pgfsetfillcolor{fc}
\pgftransformcm{1.0000}{0.0000}{0.0000}{1.0000}{\pgfqpoint{42.2654bp}{47.4427bp}}
\pgftransformscale{1.0000}
\pgftext[]{g}
\end{pgfscope}
\begin{pgfscope}
\definecolor{fc}{rgb}{0.0000,0.0000,0.0000}
\pgfsetfillcolor{fc}
\pgfsetlinewidth{0.5987bp}
\definecolor{sc}{rgb}{0.0000,0.0000,0.0000}
\pgfsetstrokecolor{sc}
\pgfsetmiterjoin
\pgfsetbuttcap
\pgfpathqmoveto{46.6057bp}{40.9871bp}
\pgfpathqcurveto{46.6057bp}{41.7698bp}{45.9712bp}{42.4044bp}{45.1884bp}{42.4044bp}
\pgfpathqcurveto{44.4057bp}{42.4044bp}{43.7711bp}{41.7698bp}{43.7711bp}{40.9871bp}
\pgfpathqcurveto{43.7711bp}{40.2043bp}{44.4057bp}{39.5698bp}{45.1884bp}{39.5698bp}
\pgfpathqcurveto{45.9712bp}{39.5698bp}{46.6057bp}{40.2043bp}{46.6057bp}{40.9871bp}
\pgfpathclose
\pgfusepathqfillstroke
\end{pgfscope}
\begin{pgfscope}
\definecolor{fc}{rgb}{0.0000,0.0000,0.0000}
\pgfsetfillcolor{fc}
\pgftransformcm{1.0000}{0.0000}{0.0000}{1.0000}{\pgfqpoint{17.4017bp}{149.6690bp}}
\pgftransformscale{1.0000}
\pgftext[]{f}
\end{pgfscope}
\begin{pgfscope}
\definecolor{fc}{rgb}{0.0000,0.0000,0.0000}
\pgfsetfillcolor{fc}
\pgfsetlinewidth{0.5987bp}
\definecolor{sc}{rgb}{0.0000,0.0000,0.0000}
\pgfsetstrokecolor{sc}
\pgfsetmiterjoin
\pgfsetbuttcap
\pgfpathqmoveto{23.3062bp}{144.1840bp}
\pgfpathqcurveto{23.3062bp}{144.9667bp}{22.6716bp}{145.6013bp}{21.8889bp}{145.6013bp}
\pgfpathqcurveto{21.1061bp}{145.6013bp}{20.4715bp}{144.9667bp}{20.4715bp}{144.1840bp}
\pgfpathqcurveto{20.4715bp}{143.4012bp}{21.1061bp}{142.7666bp}{21.8889bp}{142.7666bp}
\pgfpathqcurveto{22.6716bp}{142.7666bp}{23.3062bp}{143.4012bp}{23.3062bp}{144.1840bp}
\pgfpathclose
\pgfusepathqfillstroke
\end{pgfscope}
\begin{pgfscope}
\definecolor{fc}{rgb}{0.0000,0.0000,0.0000}
\pgfsetfillcolor{fc}
\pgftransformcm{1.0000}{0.0000}{0.0000}{1.0000}{\pgfqpoint{51.4742bp}{119.1956bp}}
\pgftransformscale{1.0000}
\pgftext[]{e}
\end{pgfscope}
\begin{pgfscope}
\definecolor{fc}{rgb}{0.0000,0.0000,0.0000}
\pgfsetfillcolor{fc}
\pgfsetlinewidth{0.5987bp}
\definecolor{sc}{rgb}{0.0000,0.0000,0.0000}
\pgfsetstrokecolor{sc}
\pgfsetmiterjoin
\pgfsetbuttcap
\pgfpathqmoveto{47.1778bp}{115.0036bp}
\pgfpathqcurveto{47.1778bp}{115.7864bp}{46.5432bp}{116.4209bp}{45.7604bp}{116.4209bp}
\pgfpathqcurveto{44.9777bp}{116.4209bp}{44.3431bp}{115.7864bp}{44.3431bp}{115.0036bp}
\pgfpathqcurveto{44.3431bp}{114.2209bp}{44.9777bp}{113.5863bp}{45.7604bp}{113.5863bp}
\pgfpathqcurveto{46.5432bp}{113.5863bp}{47.1778bp}{114.2209bp}{47.1778bp}{115.0036bp}
\pgfpathclose
\pgfusepathqfillstroke
\end{pgfscope}
\begin{pgfscope}
\definecolor{fc}{rgb}{0.0000,0.0000,0.0000}
\pgfsetfillcolor{fc}
\pgftransformcm{1.0000}{0.0000}{0.0000}{1.0000}{\pgfqpoint{72.6523bp}{74.3451bp}}
\pgftransformscale{1.0000}
\pgftext[]{d}
\end{pgfscope}
\begin{pgfscope}
\definecolor{fc}{rgb}{0.0000,0.0000,0.0000}
\pgfsetfillcolor{fc}
\pgfsetlinewidth{0.5987bp}
\definecolor{sc}{rgb}{0.0000,0.0000,0.0000}
\pgfsetstrokecolor{sc}
\pgfsetmiterjoin
\pgfsetbuttcap
\pgfpathqmoveto{70.0193bp}{80.1602bp}
\pgfpathqcurveto{70.0193bp}{80.9429bp}{69.3848bp}{81.5775bp}{68.6020bp}{81.5775bp}
\pgfpathqcurveto{67.8193bp}{81.5775bp}{67.1847bp}{80.9429bp}{67.1847bp}{80.1602bp}
\pgfpathqcurveto{67.1847bp}{79.3774bp}{67.8193bp}{78.7428bp}{68.6020bp}{78.7428bp}
\pgfpathqcurveto{69.3848bp}{78.7428bp}{70.0193bp}{79.3774bp}{70.0193bp}{80.1602bp}
\pgfpathclose
\pgfusepathqfillstroke
\end{pgfscope}
\begin{pgfscope}
\definecolor{fc}{rgb}{0.0000,0.0000,0.0000}
\pgfsetfillcolor{fc}
\pgftransformcm{1.0000}{0.0000}{0.0000}{1.0000}{\pgfqpoint{108.5059bp}{94.4302bp}}
\pgftransformscale{1.0000}
\pgftext[]{c}
\end{pgfscope}
\begin{pgfscope}
\definecolor{fc}{rgb}{0.0000,0.0000,0.0000}
\pgfsetfillcolor{fc}
\pgfsetlinewidth{0.5987bp}
\definecolor{sc}{rgb}{0.0000,0.0000,0.0000}
\pgfsetstrokecolor{sc}
\pgfsetmiterjoin
\pgfsetbuttcap
\pgfpathqmoveto{113.4253bp}{88.2694bp}
\pgfpathqcurveto{113.4253bp}{89.0521bp}{112.7907bp}{89.6867bp}{112.0079bp}{89.6867bp}
\pgfpathqcurveto{111.2252bp}{89.6867bp}{110.5906bp}{89.0521bp}{110.5906bp}{88.2694bp}
\pgfpathqcurveto{110.5906bp}{87.4866bp}{111.2252bp}{86.8520bp}{112.0079bp}{86.8520bp}
\pgfpathqcurveto{112.7907bp}{86.8520bp}{113.4253bp}{87.4866bp}{113.4253bp}{88.2694bp}
\pgfpathclose
\pgfusepathqfillstroke
\end{pgfscope}
\begin{pgfscope}
\definecolor{fc}{rgb}{0.0000,0.0000,0.0000}
\pgfsetfillcolor{fc}
\pgftransformcm{1.0000}{0.0000}{0.0000}{1.0000}{\pgfqpoint{142.1036bp}{122.4659bp}}
\pgftransformscale{1.0000}
\pgftext[]{b}
\end{pgfscope}
\begin{pgfscope}
\definecolor{fc}{rgb}{0.0000,0.0000,0.0000}
\pgfsetfillcolor{fc}
\pgfsetlinewidth{0.5987bp}
\definecolor{sc}{rgb}{0.0000,0.0000,0.0000}
\pgfsetstrokecolor{sc}
\pgfsetmiterjoin
\pgfsetbuttcap
\pgfpathqmoveto{138.8391bp}{117.1461bp}
\pgfpathqcurveto{138.8391bp}{117.9288bp}{138.2046bp}{118.5634bp}{137.4218bp}{118.5634bp}
\pgfpathqcurveto{136.6390bp}{118.5634bp}{136.0045bp}{117.9288bp}{136.0045bp}{117.1461bp}
\pgfpathqcurveto{136.0045bp}{116.3633bp}{136.6390bp}{115.7288bp}{137.4218bp}{115.7288bp}
\pgfpathqcurveto{138.2046bp}{115.7288bp}{138.8391bp}{116.3633bp}{138.8391bp}{117.1461bp}
\pgfpathclose
\pgfusepathqfillstroke
\end{pgfscope}
\begin{pgfscope}
\definecolor{fc}{rgb}{0.0000,0.0000,0.0000}
\pgfsetfillcolor{fc}
\pgftransformcm{1.0000}{0.0000}{0.0000}{1.0000}{\pgfqpoint{150.0000bp}{63.3183bp}}
\pgftransformscale{1.0000}
\pgftext[]{a}
\end{pgfscope}
\begin{pgfscope}
\definecolor{fc}{rgb}{0.0000,0.0000,0.0000}
\pgfsetfillcolor{fc}
\pgfsetlinewidth{0.5987bp}
\definecolor{sc}{rgb}{0.0000,0.0000,0.0000}
\pgfsetstrokecolor{sc}
\pgfsetmiterjoin
\pgfsetbuttcap
\pgfpathqmoveto{145.4939bp}{67.2085bp}
\pgfpathqcurveto{145.4939bp}{67.9912bp}{144.8594bp}{68.6258bp}{144.0766bp}{68.6258bp}
\pgfpathqcurveto{143.2938bp}{68.6258bp}{142.6593bp}{67.9912bp}{142.6593bp}{67.2085bp}
\pgfpathqcurveto{142.6593bp}{66.4257bp}{143.2938bp}{65.7911bp}{144.0766bp}{65.7911bp}
\pgfpathqcurveto{144.8594bp}{65.7911bp}{145.4939bp}{66.4257bp}{145.4939bp}{67.2085bp}
\pgfpathclose
\pgfusepathqfillstroke
\end{pgfscope}
\end{pgfpicture}

  \end{minipage}

  \vspace{1em}
  \begin{minipage}{0.45\textwidth}
    \begin{pgfpicture}
\pgfpathrectangle{\pgfpointorigin}{\pgfqpoint{150.0000bp}{150.0000bp}}
\pgfusepath{use as bounding box}
\begin{pgfscope}
\pgfsetlinewidth{0.5899bp}
\definecolor{sc}{rgb}{0.0000,0.0000,0.0000}
\pgfsetstrokecolor{sc}
\pgfsetmiterjoin
\pgfsetbuttcap
\pgfpathqmoveto{138.0334bp}{48.2895bp}
\pgfpathqlineto{138.0334bp}{100.0997bp}
\pgfusepathqstroke
\end{pgfscope}
\begin{pgfscope}
\pgfsetlinewidth{0.5899bp}
\definecolor{sc}{rgb}{0.0000,0.0000,0.0000}
\pgfsetstrokecolor{sc}
\pgfsetmiterjoin
\pgfsetbuttcap
\pgfpathqmoveto{138.0334bp}{100.0997bp}
\pgfpathqlineto{101.3981bp}{136.7350bp}
\pgfusepathqstroke
\end{pgfscope}
\begin{pgfscope}
\pgfsetlinewidth{0.5899bp}
\definecolor{sc}{rgb}{0.0000,0.0000,0.0000}
\pgfsetstrokecolor{sc}
\pgfsetmiterjoin
\pgfsetbuttcap
\pgfpathqmoveto{101.3981bp}{136.7350bp}
\pgfpathqlineto{138.0334bp}{48.2895bp}
\pgfusepathqstroke
\end{pgfscope}
\begin{pgfscope}
\pgfsetlinewidth{0.5899bp}
\definecolor{sc}{rgb}{0.0000,0.0000,0.0000}
\pgfsetstrokecolor{sc}
\pgfsetmiterjoin
\pgfsetbuttcap
\pgfpathqmoveto{49.5879bp}{136.7350bp}
\pgfpathqlineto{12.9525bp}{100.0997bp}
\pgfusepathqstroke
\end{pgfscope}
\begin{pgfscope}
\pgfsetlinewidth{0.5899bp}
\definecolor{sc}{rgb}{0.0000,0.0000,0.0000}
\pgfsetstrokecolor{sc}
\pgfsetmiterjoin
\pgfsetbuttcap
\pgfpathqmoveto{12.9525bp}{100.0997bp}
\pgfpathqlineto{101.3981bp}{11.6542bp}
\pgfusepathqstroke
\end{pgfscope}
\begin{pgfscope}
\pgfsetlinewidth{0.5899bp}
\definecolor{sc}{rgb}{0.0000,0.0000,0.0000}
\pgfsetstrokecolor{sc}
\pgfsetmiterjoin
\pgfsetbuttcap
\pgfpathqmoveto{12.9525bp}{48.2895bp}
\pgfpathqlineto{12.9525bp}{48.2895bp}
\pgfusepathqstroke
\end{pgfscope}
\begin{pgfscope}
\pgfsetlinewidth{0.5899bp}
\definecolor{sc}{rgb}{0.0000,0.0000,0.0000}
\pgfsetstrokecolor{sc}
\pgfsetmiterjoin
\pgfsetbuttcap
\pgfpathqmoveto{49.5879bp}{11.6542bp}
\pgfpathqlineto{49.5879bp}{11.6542bp}
\pgfusepathqstroke
\end{pgfscope}
\begin{pgfscope}
\definecolor{fc}{rgb}{0.0000,0.0000,0.0000}
\pgfsetfillcolor{fc}
\pgftransformshift{\pgfqpoint{36.6353bp}{11.6542bp}}
\pgftransformscale{1.0000}
\pgftext[]{H}
\end{pgfscope}
\begin{pgfscope}
\definecolor{fc}{rgb}{0.0000,0.0000,0.0000}
\pgfsetfillcolor{fc}
\pgfsetlinewidth{0.5899bp}
\definecolor{sc}{rgb}{0.0000,0.0000,0.0000}
\pgfsetstrokecolor{sc}
\pgfsetmiterjoin
\pgfsetbuttcap
\pgfpathqmoveto{52.1784bp}{11.6542bp}
\pgfpathqcurveto{52.1784bp}{13.0849bp}{51.0186bp}{14.2447bp}{49.5879bp}{14.2447bp}
\pgfpathqcurveto{48.1572bp}{14.2447bp}{46.9974bp}{13.0849bp}{46.9974bp}{11.6542bp}
\pgfpathqcurveto{46.9974bp}{10.2235bp}{48.1572bp}{9.0637bp}{49.5879bp}{9.0637bp}
\pgfpathqcurveto{51.0186bp}{9.0637bp}{52.1784bp}{10.2235bp}{52.1784bp}{11.6542bp}
\pgfpathclose
\pgfusepathqfillstroke
\end{pgfscope}
\begin{pgfscope}
\definecolor{fc}{rgb}{0.0000,0.0000,0.0000}
\pgfsetfillcolor{fc}
\pgftransformshift{\pgfqpoint{0.0000bp}{48.2895bp}}
\pgftransformscale{1.0000}
\pgftext[]{G}
\end{pgfscope}
\begin{pgfscope}
\definecolor{fc}{rgb}{0.0000,0.0000,0.0000}
\pgfsetfillcolor{fc}
\pgfsetlinewidth{0.5899bp}
\definecolor{sc}{rgb}{0.0000,0.0000,0.0000}
\pgfsetstrokecolor{sc}
\pgfsetmiterjoin
\pgfsetbuttcap
\pgfpathqmoveto{15.5431bp}{48.2895bp}
\pgfpathqcurveto{15.5431bp}{49.7202bp}{14.3832bp}{50.8800bp}{12.9525bp}{50.8800bp}
\pgfpathqcurveto{11.5218bp}{50.8800bp}{10.3620bp}{49.7202bp}{10.3620bp}{48.2895bp}
\pgfpathqcurveto{10.3620bp}{46.8588bp}{11.5218bp}{45.6990bp}{12.9525bp}{45.6990bp}
\pgfpathqcurveto{14.3832bp}{45.6990bp}{15.5431bp}{46.8588bp}{15.5431bp}{48.2895bp}
\pgfpathclose
\pgfusepathqfillstroke
\end{pgfscope}
\begin{pgfscope}
\definecolor{fc}{rgb}{0.0000,0.0000,0.0000}
\pgfsetfillcolor{fc}
\pgftransformshift{\pgfqpoint{110.5569bp}{2.4953bp}}
\pgftransformscale{1.0000}
\pgftext[]{F}
\end{pgfscope}
\begin{pgfscope}
\definecolor{fc}{rgb}{0.0000,0.0000,0.0000}
\pgfsetfillcolor{fc}
\pgfsetlinewidth{0.5899bp}
\definecolor{sc}{rgb}{0.0000,0.0000,0.0000}
\pgfsetstrokecolor{sc}
\pgfsetmiterjoin
\pgfsetbuttcap
\pgfpathqmoveto{103.9886bp}{11.6542bp}
\pgfpathqcurveto{103.9886bp}{13.0849bp}{102.8288bp}{14.2447bp}{101.3981bp}{14.2447bp}
\pgfpathqcurveto{99.9674bp}{14.2447bp}{98.8076bp}{13.0849bp}{98.8076bp}{11.6542bp}
\pgfpathqcurveto{98.8076bp}{10.2235bp}{99.9674bp}{9.0637bp}{101.3981bp}{9.0637bp}
\pgfpathqcurveto{102.8288bp}{9.0637bp}{103.9886bp}{10.2235bp}{103.9886bp}{11.6542bp}
\pgfpathclose
\pgfusepathqfillstroke
\end{pgfscope}
\begin{pgfscope}
\definecolor{fc}{rgb}{0.0000,0.0000,0.0000}
\pgfsetfillcolor{fc}
\pgftransformshift{\pgfqpoint{0.0000bp}{100.0997bp}}
\pgftransformscale{1.0000}
\pgftext[]{E}
\end{pgfscope}
\begin{pgfscope}
\definecolor{fc}{rgb}{0.0000,0.0000,0.0000}
\pgfsetfillcolor{fc}
\pgfsetlinewidth{0.5899bp}
\definecolor{sc}{rgb}{0.0000,0.0000,0.0000}
\pgfsetstrokecolor{sc}
\pgfsetmiterjoin
\pgfsetbuttcap
\pgfpathqmoveto{15.5431bp}{100.0997bp}
\pgfpathqcurveto{15.5431bp}{101.5304bp}{14.3832bp}{102.6902bp}{12.9525bp}{102.6902bp}
\pgfpathqcurveto{11.5218bp}{102.6902bp}{10.3620bp}{101.5304bp}{10.3620bp}{100.0997bp}
\pgfpathqcurveto{10.3620bp}{98.6690bp}{11.5218bp}{97.5092bp}{12.9525bp}{97.5092bp}
\pgfpathqcurveto{14.3832bp}{97.5092bp}{15.5431bp}{98.6690bp}{15.5431bp}{100.0997bp}
\pgfpathclose
\pgfusepathqfillstroke
\end{pgfscope}
\begin{pgfscope}
\definecolor{fc}{rgb}{0.0000,0.0000,0.0000}
\pgfsetfillcolor{fc}
\pgftransformshift{\pgfqpoint{58.7467bp}{145.8939bp}}
\pgftransformscale{1.0000}
\pgftext[]{D}
\end{pgfscope}
\begin{pgfscope}
\definecolor{fc}{rgb}{0.0000,0.0000,0.0000}
\pgfsetfillcolor{fc}
\pgfsetlinewidth{0.5899bp}
\definecolor{sc}{rgb}{0.0000,0.0000,0.0000}
\pgfsetstrokecolor{sc}
\pgfsetmiterjoin
\pgfsetbuttcap
\pgfpathqmoveto{52.1784bp}{136.7350bp}
\pgfpathqcurveto{52.1784bp}{138.1657bp}{51.0186bp}{139.3255bp}{49.5879bp}{139.3255bp}
\pgfpathqcurveto{48.1572bp}{139.3255bp}{46.9974bp}{138.1657bp}{46.9974bp}{136.7350bp}
\pgfpathqcurveto{46.9974bp}{135.3043bp}{48.1572bp}{134.1445bp}{49.5879bp}{134.1445bp}
\pgfpathqcurveto{51.0186bp}{134.1445bp}{52.1784bp}{135.3043bp}{52.1784bp}{136.7350bp}
\pgfpathclose
\pgfusepathqfillstroke
\end{pgfscope}
\begin{pgfscope}
\definecolor{fc}{rgb}{0.0000,0.0000,0.0000}
\pgfsetfillcolor{fc}
\pgftransformshift{\pgfqpoint{94.2020bp}{147.5047bp}}
\pgftransformscale{1.0000}
\pgftext[]{C}
\end{pgfscope}
\begin{pgfscope}
\definecolor{fc}{rgb}{0.0000,0.0000,0.0000}
\pgfsetfillcolor{fc}
\pgfsetlinewidth{0.5899bp}
\definecolor{sc}{rgb}{0.0000,0.0000,0.0000}
\pgfsetstrokecolor{sc}
\pgfsetmiterjoin
\pgfsetbuttcap
\pgfpathqmoveto{103.9886bp}{136.7350bp}
\pgfpathqcurveto{103.9886bp}{138.1657bp}{102.8288bp}{139.3255bp}{101.3981bp}{139.3255bp}
\pgfpathqcurveto{99.9674bp}{139.3255bp}{98.8076bp}{138.1657bp}{98.8076bp}{136.7350bp}
\pgfpathqcurveto{98.8076bp}{135.3043bp}{99.9674bp}{134.1445bp}{101.3981bp}{134.1445bp}
\pgfpathqcurveto{102.8288bp}{134.1445bp}{103.9886bp}{135.3043bp}{103.9886bp}{136.7350bp}
\pgfpathclose
\pgfusepathqfillstroke
\end{pgfscope}
\begin{pgfscope}
\definecolor{fc}{rgb}{0.0000,0.0000,0.0000}
\pgfsetfillcolor{fc}
\pgftransformshift{\pgfqpoint{150.0000bp}{105.0564bp}}
\pgftransformscale{1.0000}
\pgftext[]{B}
\end{pgfscope}
\begin{pgfscope}
\definecolor{fc}{rgb}{0.0000,0.0000,0.0000}
\pgfsetfillcolor{fc}
\pgfsetlinewidth{0.5899bp}
\definecolor{sc}{rgb}{0.0000,0.0000,0.0000}
\pgfsetstrokecolor{sc}
\pgfsetmiterjoin
\pgfsetbuttcap
\pgfpathqmoveto{140.6239bp}{100.0997bp}
\pgfpathqcurveto{140.6239bp}{101.5304bp}{139.4641bp}{102.6902bp}{138.0334bp}{102.6902bp}
\pgfpathqcurveto{136.6027bp}{102.6902bp}{135.4429bp}{101.5304bp}{135.4429bp}{100.0997bp}
\pgfpathqcurveto{135.4429bp}{98.6690bp}{136.6027bp}{97.5092bp}{138.0334bp}{97.5092bp}
\pgfpathqcurveto{139.4641bp}{97.5092bp}{140.6239bp}{98.6690bp}{140.6239bp}{100.0997bp}
\pgfpathclose
\pgfusepathqfillstroke
\end{pgfscope}
\begin{pgfscope}
\definecolor{fc}{rgb}{0.0000,0.0000,0.0000}
\pgfsetfillcolor{fc}
\pgftransformshift{\pgfqpoint{140.5603bp}{35.5858bp}}
\pgftransformscale{1.0000}
\pgftext[]{A}
\end{pgfscope}
\begin{pgfscope}
\definecolor{fc}{rgb}{0.0000,0.0000,0.0000}
\pgfsetfillcolor{fc}
\pgfsetlinewidth{0.5899bp}
\definecolor{sc}{rgb}{0.0000,0.0000,0.0000}
\pgfsetstrokecolor{sc}
\pgfsetmiterjoin
\pgfsetbuttcap
\pgfpathqmoveto{140.6239bp}{48.2895bp}
\pgfpathqcurveto{140.6239bp}{49.7202bp}{139.4641bp}{50.8800bp}{138.0334bp}{50.8800bp}
\pgfpathqcurveto{136.6027bp}{50.8800bp}{135.4429bp}{49.7202bp}{135.4429bp}{48.2895bp}
\pgfpathqcurveto{135.4429bp}{46.8588bp}{136.6027bp}{45.6990bp}{138.0334bp}{45.6990bp}
\pgfpathqcurveto{139.4641bp}{45.6990bp}{140.6239bp}{46.8588bp}{140.6239bp}{48.2895bp}
\pgfpathclose
\pgfusepathqfillstroke
\end{pgfscope}
\end{pgfpicture}

  \end{minipage}
  \begin{minipage}{0.45\textwidth}
    \begin{align*}
      G &= (V,E) \\
      V &= \{\alpha, \beta, \gamma, \delta, \epsilon, \zeta, \eta,
          \theta\} \\
      E &= \{\{\alpha, \delta\}, \{\theta, \eta\}, \{\beta, \alpha\},
          \{\zeta, \delta\}, \{\epsilon, \eta\}, \{\gamma, \alpha\}\}
    \end{align*}
    {\scriptsize
      Key: $\alpha$ = alpha, $\beta$ = beta, $\gamma$ = gamma, $\delta$ =
      delta, $\epsilon$ = epsilon, \\ $\zeta$ = zeta, $\eta$ = eta,
      $\theta$ = theta
    }
  \end{minipage}
  \end{center}

  \vspace{1em}
  \begin{defn}
    A (simple) \term{graph} $G = (V,E)$ is a set of \term{vertices}
    $V$ (singular \term{vertex}) together with a set $E$ of
    \term{edges}, where each edge consists of a set of two vertices.
  \end{defn}

  Above are shown four example graphs.  In clockwise order from the
  upper-left, we will refer to them as the numeric graph, the
  lowercase graph, the Greek graph, and the uppercase graph.  Here are
  some facts about the graphs:

  \begin{itemize}
  \item $C$ and $B$ are \term{adjacent} (aka \term{neighbors}).  So
    are $6$ and $5$.  $C$ and $F$ are not adjacent.  $f$ and $d$ are
    not adjacent.  $H$ is not adjacent to itself.
  \item The \term{degree} of vertex $C$ is 2.  The degree of vertex $H$ is
    0.  The degree of vertex $g$ is 4.
  \item Vertices $D$, $F$, $4$, $9$, $a$, $b$, $f$, $h$, $i$, $j$ are
    \term{leaves}.  No other vertex is a leaf.
  \end{itemize}
\end{model}

\begin{objective}
  Students will understand and apply graph terms \term{edge},
  \term{vertex}, \term{adjacent}/\term{neighbor}, \term{degree}, \term{leaf},
  \term{path}, \term{connected}, \term{connected component},
  \term{cycle}, \term{cyclic}, \term{acyclic}, and \term{tree}.
\end{objective}

Don't worry if you don't remember all these graph terms!  In fact, you
should probably try to forget what you think you might remember and
just focus on the information in the model. Part of the point of this
exercise is to help you either recall these terms, or learn them for
the first time.

\begin{questions}
\item How many vertices does the numeric graph have?  How many edges?
  \marginnote{\textbf{
  Important note: the singular of ``vertices'' is
  ``vertex''. ``Vertice'' is not a word.  Every time you say it, a
  puppy dies.}} \medskip
  \begin{marginfigure}
    \begin{center}
      \includegraphics[width=1in]{sad-dog}
      \caption{A sad puppy who does not want to die. Photo by Karen
        Arnold, CC$0$ Public Domain.}
    \end{center}
  \end{marginfigure}
\item How many vertices and edges does the Greek graph have?
\item Make a drawing of the Greek graph.
\item If you compared your group's drawing of the Greek graph to
  another group's drawing, do you think they would be identical, or
  different?
\item What does your answer to the previous question tell you about
  the relationship between graphs and \emph{drawings of} graphs?
\item Give another example (besides the given ones) of two vertices
  that are adjacent.
\item Give another example of two vertices that are not adjacent.
\item What does it mean for two vertices to be \term{adjacent}?
  \newpage
\item Are $\alpha$ and $\beta$ adjacent?   What about $\alpha$ and $\zeta$?
\item What is the \term{degree} of a vertex?  Use the term
  \term{adjacent} in your answer.
\item What is the degree of $\alpha$?
\item What is the definition of a \term{leaf}?  Use the term
  \term{degree} in your answer.
\item Which vertices are leaves in the Greek graph?
\item In the space below, draw a graph with at least three leaves,
  exactly one vertex with degree five, and at least one pair of
  vertices that are not adjacent.
\end{questions}

\newpage
\begin{model*}{More graph terms}{graph-terms}
Here are a few more facts about the graphs from
Model~\ref{model:graphs}:

\begin{itemize}
\item Vertices $7$--$8$--$5$--$2$--$0$ are a \term{path}.
  $2$--$1$--$0$--$2$--$1$ is also a path.  $H$ by itself is a
  path. $H$--$F$--$B$ is not a path.  $8$--$2$--$1$ is also not a path.
\item Vertices $f$ and $g$ are \term{connected}.  So are $7$ and $8$.
  $C$ and $F$ are not connected.  $H$ is connected to itself.
\item The graph with numbers is a \term{connected graph}.  So is the
  graph with lowercase letters.  The graph with uppercase letters is
  not a connected graph (it is \term{disconnected}).
\item The numeric and lowercase graphs have one \term{connected
    component} each.  The uppercase graph has four connected
  components.
\item $8$--$7$--$6$--$5$--$8$ is a \term{cycle}.  So is $C$--$A$--$B$--$C$.
  $h$--$g$--$i$--$h$ is not a cycle.  $h$--$g$--$h$ is not a cycle
  either.  Nor is $H$--$H$.  Nor is $2$--$5$--$8$--$7$--$6$--$5$.
\item The numeric graph and uppercase graph are \term{cyclic} graphs.
  The lowercase graph is \term{acyclic}.
\item The lowercase graph is a \term{tree}.  None of the other graphs
  are trees (not even the Greek one).
\end{itemize}
\end{model*}

\begin{questions}
\item How many vertices can be in a \term{path}?
\item What do you think is the definition of a \term{path}?
\item Give an example of a path in the Greek graph.
\item Can two vertices be connected but not adjacent?  If so, give an
  example.
\item Can two vertices be adjacent but not connected?  If so, give an
  example.
\item What do you think it means for two vertices to be connected?  Be
  sure to use the term \term{path} in your answer.
\item Is the Greek graph connected?
\item What do you think is the definition of a \term{connected graph}?
\item How many vertices can be in a connected component?
\item How many \term{connected components} does the Greek graph have?
\item Is the set of vertices $\{E, F\}$ a connected component?  Why or
  why not?
\item Write a definition for \term{connected component}.
\item Write an ``if and only if'' statement using the terms
  \term{connected graph} and \term{connected component}.
\item What is a \term{cycle}?  Use the term \term{path} in your
  answer.
\item Does your definition for \term{cycle} correctly explain why
  $h$--$g$--$h$ is not a cycle?  If not, revise it so it does.
  \newpage
\item Does the Greek graph have a cycle?
\item Is the Greek graph cyclic or acyclic?
\item What do you think is the definition of a tree?  You should use
  two of the other graph terms in your
  definition.\marginnote{\textbf{Warning}---a tree graph is not quite the same
    thing as a tree data structure!}
\end{questions}

\end{document}
